% Vector print adaptation of the nine published derivative-exposition figures.
% Load once; call \DerivativeFigure{key} at the published figure positions.
% Requires tikz (arrows.meta), booktabs, tabularx. No external figure files.
% Static comparison rows/panels are documented in ../figure-manifest.json.
\providecolor{ink}{HTML}{25313C}
\providecolor{inkblue}{HTML}{345B70}
\providecolor{accent}{HTML}{9B7A48}
\providecolor{wash}{HTML}{EEEAE2}
\providecolor{selection}{HTML}{E9DFCA}
\providecolor{frame}{HTML}{C8BA99}
\providecolor{pdfmuted}{HTML}{5E666B}
\providecolor{pdfline}{HTML}{DDD7CC}
\providecolor{pdfcorner}{HTML}{7A6455}
\tikzset{
  derivative text/.style={font=\fontsize{9.5}{11}\selectfont\sffamily,text=ink},
  derivative line/.style={draw=ink,line width=.8pt},
  derivative accent/.style={draw=accent,line width=1pt},
  derivative muted/.style={draw=pdfline,line width=.55pt},
  derivative arrow/.style={-{Stealth[length=4.3pt,width=3.5pt]}}
}
\newcommand{\DerivativeFigure}[1]{%
  \ifcsname derivativefigure@#1\endcsname
    \csname derivativefigure@#1\endcsname
  \else
    \PackageError{derivative-figures}{Unknown figure key '#1'}{Use one of the nine published figure keys.}%
  \fi
}
\newcommand{\DerivativeState}[1]{%
  \par\smallskip{\fontsize{9.5}{11.5}\selectfont\sffamily #1\par}\smallskip
}
\newcommand{\DerivativePowerRow}[1]{%
  \begin{tikzpicture}[x=1pt,y=1pt,derivative text]
    \path[use as bounding box] (0,0) rectangle (322,27);
    \foreach \j in {0,...,6}{%
      \pgfmathsetmacro{\px}{23+46*\j}
      \ifnum\j>#1
        \node[text=pdfmuted] at (\px,13)
          {\ifnum\j=0 \(1\)\else\ifnum\j=1 \(\varepsilon\)\else \(\varepsilon^{\j}\)\fi\fi};
        \draw[pdfmuted,line width=.55pt] (\px-9,13)--(\px+9,13);
      \else
        \fill[wash] (\px-21,0) rectangle (\px+21,27);
        \node at (\px,13)
          {\ifnum\j=0 \(1\)\else\ifnum\j=1 \(\varepsilon\)\else \(\varepsilon^{\j}\)\fi\fi};
      \fi
    }
  \end{tikzpicture}%
}

\expandafter\def\csname derivativefigure@rectangle\endcsname{%
\begin{figure}[H]
\centering
\DerivativeState{Static state: \(a=3,\ b=2,\ h=0.50\).}
\begin{tikzpicture}[x=.67pt,y=-.67pt,derivative text]
  \path[use as bounding box] (0,0) rectangle (600,355);
  \fill[wash] (90,115) rectangle (390,315);
  \fill[accent,opacity=.26] (390,115) rectangle (440,315);
  \fill[inkblue,opacity=.22] (90,65) rectangle (390,115);
  \fill[pdfcorner,opacity=.42] (390,65) rectangle (440,115);
  \draw[derivative muted] (90,115) rectangle (390,315);
  \node[anchor=base] at (240,220) {Original area \(6\)};
  \node[anchor=base] at (240,344) {\(a=3\)};
  \node[anchor=base] at (60,220) {\(b=2\)};
  \node[anchor=base] at (240,95) {\(ah\)};
  \node[anchor=base west] at (455,220) {\(bh\)};
  \node[anchor=base west] at (455,95) {\(h^2\)};
\end{tikzpicture}
\DerivativeState{Two strips \(2.50\qquad\) Corner \(0.2500\qquad\) Exact added area \(2.7500\).}
\caption{The two strips supply the linear change. The corner requires both side changes and has order two. The geometry stays visible as the increment tends to zero.}
\label{fig:rectangle}
\end{figure}%
}

\expandafter\def\csname derivativefigure@boundary\endcsname{%
\begin{figure}[H]
\centering
\DerivativeState{Static state: endpoint signs shown.}
\begin{tikzpicture}[x=.64pt,y=-.64pt,derivative text]
  \path[use as bounding box] (0,0) rectangle (600,300);
  \fill[wash] (175,60) rectangle (425,240);
  \draw[derivative line] (175,240)--(425,240)--(425,60)--(175,60)--cycle;
  \draw[derivative accent,derivative arrow] (220,240)--(380,240);
  \draw[derivative accent,derivative arrow] (425,200)--(425,105);
  \draw[derivative accent,derivative arrow] (380,60)--(220,60);
  \draw[derivative accent,derivative arrow] (175,105)--(175,200);
  \node[anchor=base] at (300,280) {bottom \(\rightarrow\)};
  \node[anchor=base] at (300,30) {\(\leftarrow\) top};
  \node[anchor=base] at (108,155) {left \(\downarrow\)};
  \node[anchor=base] at (490,155) {right \(\uparrow\)};
  \foreach \px/\py in {165/48,435/48,165/265,435/265}
    \node[anchor=base] at (\px,\py) {\(+1\;\;-1\)};
  \node[anchor=base] at (300,143) {Second boundary};
  \node[anchor=base] at (300,171) {\(0\)};
\end{tikzpicture}
\caption{Corner points cancel with opposite signs. This mechanism differs from discarding the rectangle's second-order corner area; the connection is developed through graded product rules and Stokes' theorem.}
\label{fig:boundary}
\end{figure}%
}

\expandafter\def\csname derivativefigure@diagonal\endcsname{%
\begin{figure}[H]
\centering
\begin{tikzpicture}[x=.64pt,y=-.64pt,derivative text]
  \path[use as bounding box] (0,0) rectangle (600,290);
  \draw[derivative line] (120,245)--(490,245);
  \draw[derivative line] (120,245)--(120,25);
  \node[anchor=base west] at (500,250) {\(t_{\mathrm{left}}\)};
  \node[anchor=base] at (108,18) {\(t_{\mathrm{right}}\)};
  \fill[selection] (118,221)--(433,21)--(457,59)--(142,259)--cycle;
  \draw[derivative accent] (130,240)--(445,40);
  \fill[ink] (270,151) circle[radius=3.2pt];
  \fill[inkblue] (270,128) circle[radius=3.2pt];
  \draw[derivative muted] (270,151)--(270,128);
  \node[anchor=base west] at (300,135) {difference \(h\)};
  \node[anchor=base west] at (215,205) {diagonal: \(h=0\)};
  \node[anchor=base east] at (565,65) {retain \(h\), discard \(h^2\)};
\end{tikzpicture}
\caption{The marked positions keep the left coordinate fixed; their right-coordinate difference is \(h\). Removing \(h^2\) gives the universal first-order difference. The band is schematic: an algebraic infinitesimal neighborhood has no numerical thickness.}
\label{fig:diagonal}
\end{figure}%
}

% Each arithmetic panel reuses the corresponding published circle geometry.
\newcommand{\DerivativeArithmeticPanel}[1]{%
  \begin{tikzpicture}[x=.35pt,y=-.35pt,derivative text]
    \path[use as bounding box] (0,0) rectangle (600,180);
    \ifnum#1=3
      \fill[wash] (300,75) circle[radius=11.55pt];
      \node[anchor=base] at (300,81) {\(\mathbf F_9\)};
      \node[anchor=base] at (300,145) {One closed point; two geometric points};
    \else\ifnum#1=2
      \fill[selection] (300,75) circle[radius=11.9pt];
      \fill[white] (300,75) circle[radius=7pt];
      \node[anchor=base] at (300,81) {\(1\)};
      \node[anchor=base] at (300,145) {One doubled point, length \(2\)};
    \else
      \fill[selection] (190,75) circle[radius=8.4pt];
      \fill[selection] (410,75) circle[radius=8.4pt];
      \ifnum#1=13
        \node[anchor=base] at (190,81) {\(5\)};
        \node[anchor=base] at (410,81) {\(8\)};
      \else
        \node[anchor=base] at (190,81) {\(2\)};
        \node[anchor=base] at (410,81) {\(3\)};
      \fi
      \node[anchor=base] at (300,145) {Two reduced rational points};
    \fi\fi
  \end{tikzpicture}%
}
\expandafter\def\csname derivativefigure@arithmetic\endcsname{%
\begin{figure}[H]
\centering
\DerivativeState{The online default is \(p=5\); the four available states are shown here.}
\begin{minipage}[t]{.48\linewidth}\centering
  {\fontsize{10}{12}\selectfont \(p=5:\quad\mathbf F_5\times\mathbf F_5\)}\par\smallskip
  \DerivativeArithmeticPanel{5}\par
  {\fontsize{9.5}{11.5}\selectfont Relative differential module \(0\).}
\end{minipage}\hfill
\begin{minipage}[t]{.48\linewidth}\centering
  {\fontsize{10}{12}\selectfont \(p=3:\quad\mathbf F_9\)}\par\smallskip
  \DerivativeArithmeticPanel{3}\par
  {\fontsize{9.5}{11.5}\selectfont Closed-point degree \(2\); differential module \(0\).}
\end{minipage}
\par\medskip
\begin{minipage}[t]{.48\linewidth}\centering
  {\fontsize{10}{12}\selectfont \(p=2:\quad\mathbf F_2[u]/(u^2)\)}\par\smallskip
  \DerivativeArithmeticPanel{2}\par
  {\fontsize{9.5}{11.5}\selectfont A nonzero nilpotent direction; \(\Omega\cong C\,du\).}
\end{minipage}\hfill
\begin{minipage}[t]{.48\linewidth}\centering
  {\fontsize{10}{12}\selectfont \(p=13:\quad\mathbf F_{13}\times\mathbf F_{13}\)}\par\smallskip
  \DerivativeArithmeticPanel{13}\par
  {\fontsize{9.5}{11.5}\selectfont Roots \(5,8\); relative differential module \(0\).}
\end{minipage}
\caption{The doubled point at \(2\) is detected by the differential module. The degree-two point at \(3\) becomes two points after extending the residue field.}
\label{fig:arithmetic}
\end{figure}%
}

\expandafter\def\csname derivativefigure@frobenius\endcsname{%
\begin{figure}[H]
\centering
\DerivativeState{Characteristic \(3\); input \(\varepsilon\). Online default: probe order \(m=1\).}
\DerivativePowerRow{1}
\DerivativeState{Ordinary derivative \(0\qquad\) Default output \(0\qquad\) First visible Hasse order \(3\).}
{\fontsize{10}{12}\selectfont
\begin{tabularx}{.94\linewidth}{@{}cXc@{}}
\toprule
Probe order \(m\)&Retained powers&Frobenius output\\
\midrule
\(1\)&\(1,\varepsilon\)&\(0\)\\
\(2\)&\(1,\varepsilon,\varepsilon^2\)&\(0\)\\
\(3\)&\(1,\varepsilon,\varepsilon^2,\varepsilon^3\)&\(\varepsilon^3\)\\
\bottomrule
\end{tabularx}}
\caption{For \(t\mapsto t^3\) in characteristic \(3\), an increment first survives in a probe retaining cubic terms. Increasing jet order preserves more Taylor coefficients without changing the zero first derivative.}
\label{fig:frobenius}
\end{figure}%
}

\expandafter\def\csname derivativefigure@node\endcsname{%
\begin{figure}[H]
\centering
\DerivativeState{Static default: \(a=b=1\). Residual \(ab=1.0000\): second-order lift obstructed.}
\begin{tikzpicture}[x=.64pt,y=-.64pt,derivative text]
  \path[use as bounding box] (0,0) rectangle (600,290);
  \draw[derivative muted] (115,145)--(485,145);
  \draw[derivative muted] (300,20)--(300,270);
  \draw[derivative accent] (125,145)--(475,145);
  \draw[derivative accent] (300,30)--(300,260);
  \draw[inkblue,line width=1pt,derivative arrow] (300,145)--(400,45);
  \fill[ink] (300,145) circle[radius=2.56pt];
  \fill[inkblue] (400,45) circle[radius=3.2pt];
  \node[anchor=base west] at (490,150) {\(x\)};
  \node[anchor=base] at (300,15) {\(y\)};
  \node[anchor=base west] at (320,172) {\(xy=0\)};
\end{tikzpicture}
\caption{The arrow is a tangent direction, not a point of the curve. Continuing it in the fixed node requires the product of its components to vanish.}
\label{fig:node}
\end{figure}%
}

\expandafter\def\csname derivativefigure@self-intersection\endcsname{%
\begin{figure}[H]
\centering
\begin{minipage}[t]{.48\linewidth}\centering
{\fontsize{10}{12}\selectfont\sffamily Line in \(\mathbf P^2\)}\par\smallskip
\begin{tikzpicture}[x=.66pt,y=-.66pt,derivative text]
  \path[use as bounding box] (0,0) rectangle (300,260);
  \draw[derivative line] (45,145)--(255,50);
  \draw[inkblue,line width=.85pt] (45,65)--(255,140);
  \fill[accent] (143.8235,100.2941) circle[radius=3.3pt];
  \node[anchor=base] at (150,195) {Normal bundle \(\mathcal O(1)\)};
  \node[anchor=base] at (150,228) {Self-intersection \(1\)};
\end{tikzpicture}
\end{minipage}\hfill
\begin{minipage}[t]{.48\linewidth}\centering
{\fontsize{10}{12}\selectfont\sffamily Fiber in \(\mathbf P^1\times\mathbf P^1\)}\par\smallskip
\begin{tikzpicture}[x=.66pt,y=-.66pt,derivative text]
  \path[use as bounding box] (0,0) rectangle (300,260);
  \draw[derivative line] (40,95)--(265,95);
  \draw[inkblue,line width=.85pt] (40,50)--(265,50);
  \node[anchor=base] at (150,195) {Normal bundle \(\mathcal O\)};
  \node[anchor=base] at (150,228) {Self-intersection \(0\)};
\end{tikzpicture}
\end{minipage}
\caption{A schematic displacement makes the counts visible. The intrinsic curve is \(\mathbf P^1\) in both panels; the global conormal bundle distinguishes their derived self-intersections.}
\label{fig:self-intersection}
\end{figure}%
}

\expandafter\def\csname derivativefigure@probes\endcsname{%
\begin{figure}[H]
\centering
\DerivativeState{Characteristic zero; input \(x=\varepsilon\). Online default: probe order \(m=1\).}
\DerivativePowerRow{1}
\DerivativeState{Square and cube constraints: tangent dimension \(1\), obstruction target dimension \(1\).}
{\fontsize{10}{12}\selectfont
\begin{tabularx}{.94\linewidth}{@{}cXX@{}}
\toprule
Probe order&Square constraint \(x^2=0\)&Cube constraint \(x^3=0\)\\
\midrule
\(1\)&Relation satisfied; residual \(0\)&Relation satisfied; residual \(0\)\\
\(2\)&Relation fails; residual \(\varepsilon^2\)&Relation satisfied; residual \(0\)\\
\(3\)&Relation fails; residual \(\varepsilon^2\)&Relation fails; residual \(\varepsilon^3\)\\
\bottomrule
\end{tabularx}}
\caption{The two pointwise tangent complexes agree. The square relation first fails at order \(2\), the cube relation at order \(3\). Higher brackets retain the different formal compatibility laws.}
\label{fig:probes}
\end{figure}%
}

\expandafter\def\csname derivativefigure@goodwillie\endcsname{%
\begin{figure}[H]
\centering
\begin{tikzpicture}[x=.68pt,y=-.68pt,derivative text]
  \path[use as bounding box] (0,0) rectangle (600,210);
  \filldraw[fill=wash,draw=frame,line width=.55pt] (90,100) circle[radius=30.6pt];
  \filldraw[fill=wash,draw=frame,line width=.55pt] (180,100) circle[radius=30.6pt];
  \fill[ink] (135,100) circle[radius=2.72pt];
  \node[anchor=base] at (90,106) {\(S^2\)};
  \node[anchor=base] at (180,106) {\(S^2\)};
  \node[anchor=base] at (135,180) {wedge};
  \draw[derivative line,derivative arrow] (250,100)--(345,100);
  \node[anchor=base] at (299,70) {attach \(4\)-cell};
  \node[anchor=base] at (299,140) {mixed product};
  \fill[selection] (380,48) rectangle (550,153);
  \node[anchor=base] at (465,106) {\(S^2\times S^2\)};
  \node[anchor=base] at (465,180) {product};
\end{tikzpicture}
\caption{The sphere outlines and product block are schematic. The relevant datum is the attaching map of the 4-cell, the mixed Whitehead product. The first stable approximation loses this interaction; the second derivative detects it.}
\label{fig:goodwillie}
\end{figure}%
}
